\documentclass[11pt,a4paper]{article}
\usepackage[utf8]{inputenc}
\usepackage[spanish]{babel}
\usepackage{geometry}
\geometry{top=2.5cm, bottom=2.5cm, left=2.5cm, right=2.5cm}
\usepackage{graphicx}
\usepackage{listings}
\usepackage{xcolor}
\usepackage{amsmath}
\usepackage{amsfonts}
\usepackage{booktabs}
\usepackage{hyperref}

\definecolor{codegreen}{rgb}{0,0.6,0}
\definecolor{codegray}{rgb}{0.5,0.5,0.5}
\definecolor{codeblue}{rgb}{0,0,0.8}
\definecolor{codepurple}{rgb}{0.58,0,0.82}
\definecolor{backcolour}{rgb}{0.96,0.96,0.98}

\lstdefinestyle{mystyle}{
    backgroundcolor=\color{backcolour},   
    commentstyle=\color{codegreen},
    keywordstyle=\color{codeblue}\bfseries,
    numberstyle=\tiny\color{codegray},
    stringstyle=\color{codepurple},
    basicstyle=\ttfamily\small,
    breakatwhitespace=false,         
    breaklines=true,                 
    captionpos=b,                    
    keepspaces=true,                 
    numbers=left,                    
    numbersep=5pt,                  
    showspaces=false,                
    showstringspaces=false,
    showtabs=false,                  
    tabsize=4,
    language=C++
}
\lstset{style=mystyle}

\title{
    \textbf{Universidad Nacional de San Agustín de Arequipa} \\
    \textbf{Escuela Profesional de Ciencia de la Computación} \\
    \vspace{0.5cm}
    \includegraphics[width=0.3\textwidth]{logo_unsa.png} \\ % Reemplazar con el logo correspondiente
    \vspace{0.5cm}
    \large{Computación Gráfica - Laboratorio 10} \\
    \Huge{\textbf{Simulación de Movimiento Colectivo: Boids}}
}
\author{\textbf{Estudiante}: low}
\date{\today}

\begin{document}

\maketitle
\newpage
\tableofcontents
\newpage

\section{Introducción}
El modelo de Boids, propuesto por Craig Reynolds en 1986, es un ejemplo clásico de comportamiento emergente, donde reglas individuales simples dan lugar a un comportamiento colectivo complejo y coordinado. La simulación modela el vuelo en bandada de aves u otras criaturas gregarias utilizando tres reglas fundamentales de dirección basadas en la vecindad local: Separación, Alineamiento y Cohesión.

En este informe se detalla la implementación en C++ utilizando OpenGL y GLUT de una simulación interactiva de Boids, incluyendo un panel de control interactivo que permite modificar todos los parámetros físicos y de la población en tiempo real.

\newpage

\section{Desarrollo del Cuestionario de Laboratorio}

\subsection{Punto 1: Estructura del Boid}
\textbf{Descripción}: \\
Para representar a cada agente individual (Boid) en el espacio bidimensional de la simulación, se implementa una estructura estructurada en C++ que almacena de forma eficiente su estado de posición ($x, y$), su velocidad de movimiento en ambos ejes ($vx, vy$), su orientación angular o dirección, su identificador único para aplicar un color distintivo, y un historial circular de posiciones anteriores para la generación y renderizado de estelas de movimiento.

\textbf{Fracción de Código}:
\begin{lstlisting}[caption={Definición de la estructura Boid en boid.h}, label={lst:boid_struct}]
struct Boid {
    float x, y;
    float vx, vy;
    float direction;
    int id;

    float trailX[MAX_TRAIL];
    float trailY[MAX_TRAIL];
    int trailCount;

    void updateDirection() {
        direction = atan2(vy, vx);
    }

    void savePosition() {
        for (int i = MAX_TRAIL - 1; i > 0; i--) {
            trailX[i] = trailX[i - 1];
            trailY[i] = trailY[i - 1];
        }
        trailX[0] = x;
        trailY[0] = y;
        if (trailCount < MAX_TRAIL) trailCount++;
    }
};
\end{lstlisting}

\textbf{Explicación del Código}: \\
La estructura \lstinline{Boid} encapsula la posición flotante, el vector velocidad, y la dirección calculada a través de la función arcotangente \lstinline{updateDirection()}. Adicionalmente, el método \lstinline{savePosition()} actúa como un búfer de desplazamiento temporal para desplazar las posiciones previas y almacenar la posición actual en el índice $0$, permitiendo renderizar una estela que se difumina progresivamente en el tiempo sin requerir memoria dinámica adicional.

\textbf{Evidencia}:
\begin{figure}[htbp]
    \centering
    % \includegraphics[width=0.8\textwidth]{evidencia_punto1.png}
    \caption{Representación en ejecución de un Boid individual y su historial de movimiento.}
    \label{fig:evidencia_punto1}
\end{figure}

\newpage

\subsection{Punto 2: Inicialización de la Población}
\textbf{Descripción}: \\
La población de boids se inicializa dinámicamente según la cantidad configurada en la UI. Para garantizar que los agentes no aparezcan detrás del panel lateral de control cuando este se encuentra visible al inicio, la región de generación aleatoria se restringe al espacio visible de la simulación. A cada agente se le asigna una dirección angular inicial aleatoria en el rango $[0, 2\pi]$ y una magnitud de velocidad inicial uniforme acotada entre los límites configurados.

\textbf{Fracción de Código}:
\begin{lstlisting}[caption={Inicialización aleatoria de Boids en simulacion.cpp}, label={lst:init_boids}]
void initializeBoids(int n) {
    boids.clear();
    float minX = uiPanelOpen ? PANEL_W : 0.0f;
    float simW = static_cast<float>(WINDOW_W) - minX;

    for (int i = 0; i < n; i++) {
        Boid b;
        b.x = minX + 50.0f + static_cast<float>(rand() % static_cast<int>(simW - 100.0f));
        b.y = 50.0f + static_cast<float>(rand() % (WINDOW_H - 100));
        b.id = i;
        b.trailCount = 0;

        for (int t = 0; t < MAX_TRAIL; t++) {
            b.trailX[t] = b.x;
            b.trailY[t] = b.y;
        }

        float angulo = static_cast<float>(rand()) / RAND_MAX * 2.0f * M_PI;
        float velocidad = minSpeed + static_cast<float>(rand()) / RAND_MAX * (maxSpeed - minSpeed);
        b.vx = cos(angulo) * velocidad;
        b.vy = sin(angulo) * velocidad;
        b.updateDirection();

        boids.push_back(b);
    }
}
\end{lstlisting}

\textbf{Explicación del Código}: \\
El método calcula el ancho de pantalla utilizable (\lstinline{simW}) restando el ancho del panel si este está abierto. Luego distribuye los $n$ boids mediante números pseudoaleatorios en coordenadas de pantalla válidas y define sus componentes de velocidad $vx, vy$ multiplicando el seno y coseno del ángulo aleatorio por una velocidad aleatoria comprendida entre \lstinline{minSpeed} y \lstinline{maxSpeed}.

\textbf{Evidencia}:
\begin{figure}[htbp]
    \centering
    % \includegraphics[width=0.8\textwidth]{evidencia_punto2.png}
    \caption{Población inicial de boids distribuida uniformemente fuera de la zona del panel de control.}
    \label{fig:evidencia_punto2}
\end{figure}

\newpage

\subsection{Punto 3: Representación Visual Orientada}
\textbf{Descripción}: \\
Para reemplazar los tradicionales triángulos estáticos por una representación atractiva y dinámica, cada boid es dibujado como una silueta estilizada de ave en 2D. La orientación visual del ave coincide de forma exacta con su dirección de movimiento actual mediante la transformación de rotación de la matriz modelo-vista en OpenGL basada en la variable \lstinline{direction}. Además, se añade una animación de aleteo en las alas, cuya frecuencia y amplitud dependen de la velocidad de traslación del boid.

\textbf{Fracción de Código}:
\begin{lstlisting}[caption={Renderizado y orientación del boid en render.cpp}, label={lst:draw_boid}]
void drawBoid(const Boid& b) {
    float hue = fmod(static_cast<float>(b.id) * 0.073f, 1.0f);
    float r, g, bl;
    hsvToRgb(hue, 0.8f, 1.0f, r, g, bl);

    float rOscuro, gOscuro, blOscuro;
    hsvToRgb(hue, 0.9f, 0.5f, rOscuro, gOscuro, blOscuro);

    float speed = sqrt(b.vx * b.vx + b.vy * b.vy);
    float phase = globalTime * 8.0f + b.id * 1.3f;
    float flap = sin(phase) * (0.3f + speed * 0.1f);

    float S = 1.0f;

    glPushMatrix();
    glTranslatef(b.x, b.y, 0.0f);
    glRotatef(b.direction * 180.0f / M_PI, 0.0f, 0.0f, 1.0f);

    // Renderizado del cuerpo, cabeza, pico y alas animadas usando primitivas GL_TRIANGLES y GL_POLYGON...
    // ...
    glPopMatrix();
}
\end{lstlisting}

\textbf{Explicación del Código}: \\
Se traslada la matriz de transformación a la posición del boid ($b.x, b.y$) y se aplica una rotación angular \lstinline{glRotatef} usando el ángulo de movimiento convertido a grados. Para simular el aleteo, las coordenadas de los vértices de las alas se desplazan multiplicándolas por la variable armónica \lstinline{flap}. Se realiza una conversión HSV a RGB única por ID para asegurar que cada boid tenga una tonalidad de color diferenciada y estéticamente atractiva.

\textbf{Evidencia}:
\begin{figure}[htbp]
    \centering
    % \includegraphics[width=0.8\textwidth]{evidencia_punto3.png}
    \caption{Detalle de las aves orientadas en la dirección de su velocidad con su respectivo aleteo y color asignado.}
    \label{fig:evidencia_punto3}
\end{figure}

\newpage

\subsection{Punto 4: Movimiento Continuo y Reglas de Reynolds}
\textbf{Descripción}: \\
Los boids se mueven continuamente por el lienzo actualizando sus coordenadas según su velocidad acumulada. Para implementar el comportamiento colectivo, en cada fotograma se evalúan las tres reglas de Reynolds:
\begin{itemize}
    \item \textbf{Separación}: Fuerza de repulsión inversa a la distancia respecto a boids muy cercanos para evitar colisiones físicas directas.
    \item \textbf{Alineamiento}: Fuerza que empuja al boid a orientar su vector velocidad en la misma dirección que el promedio de velocidades de sus vecinos locales.
    \item \textbf{Cohesión}: Fuerza de atracción hacia el centro de masa promedio de los boids vecinos para mantener unida a la bandada.
\end{itemize}

\textbf{Fracción de Código}:
\begin{lstlisting}[caption={Cálculo de las fuerzas físicas en simulacion.cpp}, label={lst:reynolds_rules}]
void ruleSeparation(int i, float& fx, float& fy) {
    fx = 0.0f; fy = 0.0f;
    int neighbors = 0;
    for (int j = 0; j < static_cast<int>(boids.size()); j++) {
        if (i == j) continue;
        float dx, dy, dist;
        calculateDistance(boids[i], boids[j], dx, dy, dist);
        if (dist < separationDistance && dist > 0.001f) {
            fx -= dx / dist;
            fy -= dy / dist;
            neighbors++;
        }
    }
    if (neighbors > 0) { fx /= neighbors; fy /= neighbors; }
}

void ruleAlignment(int i, float& fx, float& fy) {
    fx = 0.0f; fy = 0.0f;
    float avgVx = 0.0f, avgVy = 0.0f;
    int neighbors = 0;
    for (int j = 0; j < static_cast<int>(boids.size()); j++) {
        if (i == j) continue;
        float dx, dy, dist;
        calculateDistance(boids[i], boids[j], dx, dy, dist);
        if (dist < neighborRadius) {
            avgVx += boids[j].vx;
            avgVy += boids[j].vy;
            neighbors++;
        }
    }
    if (neighbors > 0) {
        avgVx /= neighbors; avgVy /= neighbors;
        fx = avgVx - boids[i].vx; fy = avgVy - boids[i].vy;
    }
}
\end{lstlisting}

\textbf{Explicación del Código}: \\
Cada regla itera sobre la población de boids, descarta al propio boid y mide la distancia. Si la distancia está dentro de los límites requeridos (\lstinline{separationDistance} para repulsión y \lstinline{neighborRadius} para alineamiento y cohesión), acumula las contribuciones relativas. Al final, se divide por el número de vecinos detectados para obtener vectores promedio.

\textbf{Evidencia}:
\begin{figure}[htbp]
    \centering
    % \includegraphics[width=0.8\textwidth]{evidencia_punto4.png}
    \caption{Visualización de agrupaciones coordinadas que surgen al interactuar las tres fuerzas dinámicas.}
    \label{fig:evidencia_punto4}
\end{figure}

\newpage

\subsection{Punto 5: Combinación y Ponderación de Reglas}
\textbf{Descripción}: \\
Las tres fuerzas vectoriales resultantes se multiplican por sus pesos de ponderación correspondientes y se aplican a la aceleración del boid, modificando su velocidad acumulada en cada fotograma.

\textbf{Fracción de Código}:
\begin{lstlisting}[caption={Combinación lineal de fuerzas ponderadas en simulacion.cpp}, label={lst:combine_rules}]
void updateBoids() {
    if (simulationPaused) return;
    int n = static_cast<int>(boids.size());

    for (int i = 0; i < n; i++) {
        float sepX, sepY;
        ruleSeparation(i, sepX, sepY);

        float aliX, aliY;
        ruleAlignment(i, aliX, aliY);

        float cohX, cohY;
        ruleCohesion(i, cohX, cohY);

        boids[i].vx += sepX * separationWeight
                     + aliX * alignmentWeight * 0.05f
                     + cohX * cohesionWeight  * 0.05f;

        boids[i].vy += sepY * separationWeight
                     + aliY * alignmentWeight * 0.05f
                     + cohX * cohesionWeight  * 0.05f;

        limitSpeed(boids[i]);
        applyBoundaries(boids[i]);

        boids[i].savePosition();
        boids[i].x += boids[i].vx;
        boids[i].y += boids[i].vy;
        boids[i].updateDirection();
    }
    globalTime += 0.05f;
}
\end{lstlisting}

\textbf{Explicación del Código}: \\
El integrador de Euler actualiza la velocidad sumándole las fuerzas acumuladas ponderadas por \lstinline{separationWeight}, \lstinline{alignmentWeight} y \lstinline{cohesionWeight}. El factor multiplicador $0.05$ se utiliza para suavizar la influencia de la cohesión y alineamiento, de modo que las maniobras colectivas se realicen de forma fluida sin cambios bruscos de dirección.

\textbf{Evidencia}:
\begin{figure}[htbp]
    \centering
    % \includegraphics[width=0.8\textwidth]{evidencia_punto5.png}
    \caption{Control y cambio en tiempo real del comportamiento al alternar las ponderaciones desde el menú lateral.}
    \label{fig:evidencia_punto5}
\end{figure}

\newpage

\subsection{Punto 6: Límites de Velocidad}
\textbf{Descripción}: \\
Para evitar que las fuerzas acumuladas aceleren a los boids indefinidamente (haciéndolos salir del rango visual instantáneamente) o los dejen estáticos en la pantalla en ausencia de fuerzas vecinas, se implementa una función reguladora que restringe la velocidad del boid manteniéndola siempre en un intervalo acotado de velocidad mínima y máxima.

\textbf{Fracción de Código}:
\begin{lstlisting}[caption={Ajuste de límites de velocidad en simulacion.cpp}, label={lst:limit_speed}]
void limitSpeed(Boid& b) {
    float speed = sqrt(b.vx * b.vx + b.vy * b.vy);
    if (speed < 0.0001f) return;

    if (speed > maxSpeed) {
        b.vx = (b.vx / speed) * maxSpeed;
        b.vy = (b.vy / speed) * maxSpeed;
    } else if (speed < minSpeed) {
        b.vx = (b.vx / speed) * minSpeed;
        b.vy = (b.vy / speed) * minSpeed;
    }
}
\end{lstlisting}

\textbf{Explicación del Código}: \\
Calculamos la magnitud de la velocidad mediante la hipotenusa de los vectores. Si la velocidad excede \lstinline{maxSpeed}, normalizamos el vector dividiéndolo por su magnitud actual y multiplicándolo por la velocidad máxima permitida. De igual forma, si es menor a \lstinline{minSpeed}, forzamos a que tenga al menos esa magnitud para mantener el movimiento constante de la bandada.

\textbf{Evidencia}:
\begin{figure}[htbp]
    \centering
    % \includegraphics[width=0.8\textwidth]{evidencia_punto6.png}
    \caption{Movimiento fluido donde las velocidades se ven limitadas de acuerdo al rango del deslizador doble.}
    \label{fig:evidencia_punto6}
\end{figure}

\newpage

\subsection{Punto 7: Comportamiento en los Bordes de la Ventana}
\textbf{Descripción}: \\
Para evitar que los agentes abandonen permanentemente el área simulación, se han implementado dos modos de comportamiento de frontera seleccionables dinámicamente:
\begin{itemize}
    \item \textbf{Modo Rebote}: Una fuerza virtual invisible desvía al boid a medida que se acerca al borde del panel o de la ventana. En caso de traspaso crítico (debido a fuerzas acumuladas intensas), se realiza una corrección instantánea de posición y rebote elástico invirtiendo la velocidad.
    \item \textbf{Mundo Toroidal}: El boid se teletransporta al borde opuesto al traspasar la frontera visible. Para evitar estelas indeseadas cruzando la pantalla de lado a lado, se añadió un filtro de distancia que suprime el dibujo de estelas que exceden un umbral de distancia.
\end{itemize}

\textbf{Fracción de Código}:
\begin{lstlisting}[caption={Control de límites espaciales en simulacion.cpp}, label={lst:boundaries}]
void applyBoundaries(Boid& b) {
    float minX = uiPanelOpen ? PANEL_W : 0.0f;
    float maxX = static_cast<float>(WINDOW_W);
    float minY = 0.0f;
    float maxY = static_cast<float>(WINDOW_H);

    if (edgeMode == 0) { // Rebote
        if (b.x < minX + edgeMargin) b.vx += edgeForce;
        if (b.x > maxX - edgeMargin) b.vx -= edgeForce;
        if (b.y < minY + edgeMargin) b.vy += edgeForce;
        if (b.y > maxY - edgeMargin) b.vy -= edgeForce;

        if (b.x < minX) { b.x = minX; b.vx = std::abs(b.vx); }
        else if (b.x > maxX) { b.x = maxX; b.vx = -std::abs(b.vx); }
        if (b.y < minY) { b.y = minY; b.vy = std::abs(b.vy); }
        else if (b.y > maxY) { b.y = maxY; b.vy = -std::abs(b.vy); }
    } else { // Toroidal
        if (b.x < minX)         b.x += (maxX - minX);
        if (b.x > maxX)         b.x -= (maxX - minX);
        if (b.y < minY)         b.y += (maxY - minY);
        if (b.y > maxY)         b.y -= (maxY - minY);
    }
}
\end{lstlisting}

\textbf{Explicación del Código}: \\
Si \lstinline{edgeMode == 0}, se calculan márgenes virtuales (\lstinline{edgeMargin}) y se aplica una aceleración contraria (\lstinline{edgeForce}) al aproximarse al borde para una transición suave. Adicionalmente, el chequeo duro asigna valores absolutos apropiados si el agente sobrepasa físicamente la frontera de la simulación. En el modo toroidal, se aplica aritmética de envoltura espacial para reposicionar al boid al otro lado de la pantalla.

\textbf{Evidencia}:
\begin{figure}[htbp]
    \centering
    % \includegraphics[width=0.8\textwidth]{evidencia_punto7.png}
    \caption{Agentes rodeando los límites del panel dinámico (rebote) o reapareciendo en el extremo opuesto (toroidal).}
    \label{fig:evidencia_punto7}
\end{figure}

\newpage

\subsection{Punto 8: Modificaciones de Parámetros en Tiempo Real}
\textbf{Descripción}: \\
Se ha diseñado un panel de control avanzado que se integra dentro de la simulación mediante un sidebar translúcido y colapsable. Este panel permite modificar de forma inmediata la cantidad de boids, el radio de vecindad local, las constantes de ponderación de las reglas físicas, los modos de borde (Rebote o Toroidal), visualizar y ocultar las estelas de movimiento mediante un Checkbox y pausar/reanudar el flujo temporal. La posición del ratón sobre la ventana física de GLUT se mapea linealmente y escala en tiempo real para soportar cualquier aspecto de pantalla de administradores de ventanas como Hyprland.

\textbf{Fracción de Código}:
\begin{lstlisting}[caption={Coordenadas escaladas para el manejo del ratón en ui.cpp}, label={lst:mouse_scaling}]
bool handleMouse(int button, int state, int x, int y) {
    if (button != GLUT_LEFT_BUTTON) return false;

    float actualW = static_cast<float>(glutGet(GLUT_WINDOW_WIDTH));
    float actualH = static_cast<float>(glutGet(GLUT_WINDOW_HEIGHT));
    if (actualW < 1.0f) actualW = 1.0f;
    if (actualH < 1.0f) actualH = 1.0f;

    float mx = static_cast<float>(x) * (static_cast<float>(WINDOW_W) / actualW);
    float my = (actualH - static_cast<float>(y)) * (static_cast<float>(WINDOW_H) / actualH);
    
    // Evaluacion de clics sobre botones, sliders, checkbox y controles segmentados...
    // ...
    return true;
}
\end{lstlisting}

\textbf{Explicación del Código}: \\
Para solucionar el desfase de la posición del ratón en ventanas redimensionadas bajo Hyprland, \lstinline{glutGet} consulta dinámicamente las dimensiones físicas de la ventana (\lstinline{actualW, actualH}). Las posiciones de la UI se mapean multiplicando la coordenada del ratón por la relación de la proyección virtual ($1000 \times 700$) y la altura invertida de OpenGL, permitiendo una precisión de pixel perfecta independientemente del tamaño de pantalla del usuario.

\textbf{Evidencia}:
\begin{figure}[htbp]
    \centering
    % \includegraphics[width=0.8\textwidth]{evidencia_punto8.png}
    \caption{Interfaz lateral de usuario modificando dinámicamente los parámetros de la simulación.}
    \label{fig:evidencia_punto8}
\end{figure}

\newpage

\section{Conclusiones}
\begin{itemize}
    \item La arquitectura de Boids basada en vecindad local demuestra con claridad que sistemas no jerárquicos y descentralizados pueden lograr estados globales coordinados estables, similares a los observados en bandadas reales de aves.
    \item El modelado físico combinando la repulsión elástica progresiva en los bordes y la envoltura toroidal permite probar distintas geometrías dinámicas de fluidos y comportamientos biológicos virtuales con precisión matemática.
    \item El uso de una proyección ortográfica estática combinada con escalado dinámico de eventos de ratón en tiempo de ejecución es fundamental para lograr una interfaz interactiva robusta y compatible con gestores de ventanas dinámicos modernos en Linux.
\end{itemize}

\end{document}
